%!TeX root = ../main.tex
% ---------------------------------------------------------------
\section{Prospects}
\label{sec:prospects}
% ---------------------------------------------------------------

The paper establishes that ordinal spectral ranking of journals and institutions
by recirculation is feasible, interpretable against expert rankings, and largely stable
across a quarter-century of economics and business literature. Several threads
invite further development.

A useful extension could be work-level spectral ranking over the directed
acyclic graph of work-work references. As noted in Section~\ref{sec:discussion},
projecting works to units introduces fictitious information by conflating independent
reference chains passing through a shared unit. Spectral methods applied directly to
the work-level graph could avoid this conflation, but currently call on the Katz-Hubbell 
(PageRank) resolvent. While $\alpha$ and $\mathbf{\mu}$ stabilise the solution they 
distort the network, defeating the intent. High order eigenpairs would add information 
that could help resolve the trade-off.

The sentinel analysis ($\varepsilon=1$) identifies units whose reference flows are
connected to works outside the corpus boundary. Broadening the network to a more 
faithful representation will generally increase the presence of sub-structure in the network, 
potentially clouding the interpretation of the wider ranking. A systematic
characterisation of these boundary flows would complement unit-level influence scores
and distinguish units whose standing is field-specific from those anchored across a
broader literature. Again, high order eigenpairs could help delineate these tensions.

The enclave phenomenon uncovers communities of highly cited works in low-influence
units. As flagged in the Introduction, actor-network theory could be important for
understanding the intent of researchers, editors and managers within these communities
\citep{latourReassemblingSocialIntroduction2007}. Connecting enclave membership to
authorship networks, career trajectories and funding structures is a natural follow-on
study, as is investigating whether other fields also have enclaves.

Finally, extending the corpus to additional disciplines and adapting the methodology
for performance-based research funding schemes will require field-specific adjustments
to the source registry and corpus pipeline, but the framework is otherwise directly
applicable.

